\subsection{Problema inverso en radioterapia}

Una vez introducidas las herramientas numéricas necesarias para resolver problemas de optimización multivariable, resulta posible formular específicamente el problema de planificación inversa en radioterapia IMRT.

La radioterapia moderna busca administrar una dosis suficientemente elevada de radiación sobre el volumen tumoral, maximizando la probabilidad de control del tumor, mientras se minimiza simultáneamente la irradiación recibida por tejidos sanos y órganos críticos circundantes. Este objetivo introduce naturalmente un problema de compromiso entre múltiples requerimientos físicos y clínicos que, en general, resultan incompatibles entre sí.

En técnicas convencionales de planificación directa, el operador selecciona manualmente parámetros geométricos del tratamiento —como direcciones de incidencia, formas de campo e intensidades de irradiación— y posteriormente evalúa la distribución de dosis resultante. Este procedimiento depende fuertemente de la experiencia del planificador y suele requerir múltiples iteraciones manuales para alcanzar resultados clínicamente aceptables.

La radioterapia de intensidad modulada (\textit{Intensity-Modulated Radiation Therapy}, IMRT) reformula este proceso mediante un enfoque de planificación inversa. En lugar de especificar directamente las intensidades de irradiación, el clínico establece previamente objetivos dosimétricos deseados sobre distintas regiones anatómicas. El problema consiste entonces en determinar automáticamente la configuración de intensidades capaz de producir una distribución de dosis lo más cercana posible a dichos objetivos.

Desde el punto de vista matemático, el problema recibe el nombre de problema inverso debido a que las causas físicas del sistema —las intensidades de irradiación— no son conocidas inicialmente, mientras que sí se conocen los efectos deseados sobre la distribución de dosis. El objetivo consiste, por lo tanto, en reconstruir las variables de control del sistema a partir de restricciones impuestas sobre su respuesta.

En IMRT, cada haz de radiación se discretiza en un conjunto de sub-haces elementales denominados \textit{beamlets}. Cada beamlet aporta una determinada cantidad de dosis sobre los distintos vóxeles del paciente. Debido a que la deposición total de dosis puede aproximarse mediante superposición lineal de contribuciones individuales, la dosis absorbida en cada vóxel puede expresarse como combinación lineal de las intensidades asociadas a todos los beamlets.

Esta discretización transforma el problema físico continuo en un problema de optimización multivariable de gran dimensión. Las variables desconocidas corresponden a las intensidades de los beamlets, mientras que las restricciones clínicas se traducen en objetivos dosimétricos sobre regiones anatómicas específicas.

Sin embargo, las exigencias clínicas del tratamiento suelen ser inherentemente conflictivas. Incrementar la dosis administrada al volumen tumoral puede producir simultáneamente un aumento indeseado de dosis sobre órganos de riesgo cercanos. Del mismo modo, reducir excesivamente la irradiación sobre tejido sano puede comprometer la cobertura tumoral requerida. Como consecuencia, el problema no posee generalmente una solución exacta capaz de satisfacer simultáneamente todos los objetivos clínicos impuestos.

En este contexto, la planificación inversa se formula como un problema de optimización numérica cuyo objetivo consiste en encontrar una configuración de intensidades que minimice globalmente las desviaciones respecto de los criterios clínicos deseados. La complejidad matemática del problema surge tanto de la elevada dimensionalidad del espacio de variables como de la necesidad de equilibrar simultáneamente múltiples objetivos competitivos.

Debido a estas características, la resolución eficiente del problema requiere la utilización de algoritmos de optimización capaces de operar sobre funciones objetivo de gran dimensión con costos computacionales razonables. Entre los métodos más utilizados para este tipo de aplicaciones se encuentran los algoritmos quasi-Newton de memoria limitada, particularmente L-BFGS, debido a su capacidad para incorporar información aproximada de curvatura sin necesidad de almacenar explícitamente matrices Hessianas de gran tamaño.